% Appendix B -- Model details (draft prose; from the R0 code)
%   budgets      estimation_epc/model_funcs_estimation_epc.py: outcomes(),
%                _buy_cost, sell_proceeds_t, _owner_recurring_cost,
%                _stay_payment_and_balance, _buy_effective_n_rem, _stay_n_rem_next
%   debt caps    epc_revised/economics.py: buy_terms, stay_min_payment,
%                permanent_income
%   82+ group    epc_revised/spec.py (T = 60, recurring_terminal = 1),
%                simulation horizon 70 periods
\section{Model Details}\label{app:model}

\subsection{The 82+ group and the simulation.}\label{app:model-terminal}
The model has 60 ages, $J_0=23$ to $\bar J=82$. The last stands for all
ages from $\bar J$ on: a household in this group survives each year with
probability $\phi_{\bar J}$, and its income, equivalence scale, and choice
set do not change. Its value function is therefore the same every year and
solves a fixed point. Simulated households are followed for 70 years from
age $J_0$; those still alive after $\bar J$ follow the policy of age $\bar J$.

\subsection{Mortgages.}\label{app:model-mortgages}
Let $t_0$ be the efficiency of an owner's current home and
$B=(1+r_m)m-P(m,n)$ the balance due next year under the existing
schedule. The rate $r_m$ is the same for all loans. The remaining term falls by one
each year. A loan taken by a household without a mortgage has term $L$. A
mover with a positive balance keeps its remaining term, which falls by one
as for a stayer. When one year remains, the payment
$P(m,1)=(1+r_m)m$ repays the balance. The repayment schedule continues
in retirement; there is no interest-only phase.

\emph{Purchase.} A buyer's loan satisfies
\begin{equation*}
  m'\leq\min\bigl\{\lambda_Lp^th',\ P^{-1}\bigl(\lambda_Dy^{\mathrm{perm}};\,n'\bigr)\bigr\},
\end{equation*}
where $P^{-1}(\cdot;n')$ is the largest balance whose payment over $n'$
years equals its argument, $n'$ is the term of the new loan, and $y^{\mathrm{perm}}$ is
gross income with the transitory shock set to one (the pension in
retirement). For a purchase with renovation, $t=N$. An owner who moves
also faces
\begin{equation*}
  m'\leq B+\max\bigl\{p^th'-p^{t_0}h,\ 0\bigr\}+C(h')\,\ind\{\text{renovates at purchase}\}.
\end{equation*}
A mover can therefore borrow for a more valuable home and for a renovation
at purchase, but cannot take equity out; a move to a home of equal value
continues the existing schedule.

\emph{Renovation in place.} An owner who renovates in place receives the
advance
\begin{equation*}
  A=\min\bigl\{C(h),\ \max\{K-B,\ 0\}\bigr\},
  \qquad
  K=\min\bigl\{\lambda_Lp^Nh,\ P^{-1}\bigl(\lambda_Dy^{\mathrm{perm}};\,\hat n\bigr)\bigr\},
\end{equation*}
where $\hat n=n-1$ for an owner with a mortgage and $\hat n=L$ for an owner
without one. The advance is added to next year's balance and is repaid by
the existing payoff date. The collateral is the inefficient value of the
home.

\subsection{Initial conditions.}\label{app:model-initial}
Households enter at age 23 with the income, liquid wealth, tenure, home
value, mortgage balance, and remaining term of records of household heads
aged 20--29 in the Wealth and Assets Survey; Appendix~\ref{app:data-was}
describes the sample. Records are drawn in proportion to their weights.
Income in the record sets the permanent component, with a floor, and the
transitory component starts at one. Negative net financial wealth is set
to zero, since the model has no unsecured debt. An owner's home size is its
value in the record, kept within the limits $[\underline h^{\mathrm{own}},\bar
h^{\mathrm{own}}]$. An owner is placed in an efficient home with
probability $\pi^R_0$; the market value of its home is kept, so the size
of an efficient home is its value divided by $p^R$.

\subsection{Budget constraints.}\label{app:model-budgets}
Let $x$ be cash on hand. Consumption is the residual in each budget
constraint below.

\emph{Renting.} A renter who keeps renting, or an owner who sells and
rents, pays
\begin{equation*}
  c+b'+\bigl[p^{\mathrm{rent}}+\pi_Re^R+(1-\pi_R)e^N\bigr]h'=x+S,
\end{equation*}
where $S=0$ for a renter and, for an owner who sells, $S$ is the sale
proceeds
\begin{equation*}
  S=(1-\delta_h-\zeta_s)\,p^{t_0}h-(1+r_m)m.
\end{equation*}
The seller pays the year's maintenance on the home it sells and repays the
mortgage in full. Rental size lies between $\underline h^{\mathrm{rent}}$ and
$\bar h^{\mathrm{rent}}$.

\emph{Buying.} A household that buys a home of size $h'$ and state $t$
pays
\begin{equation*}
  c+b'+(1+\zeta_b)\,p^th'+e^th'=x+S+m'.
\end{equation*}
A household that buys an inefficient home and renovates it pays
$(1+\zeta_b)p^Nh'+C(h')+e^Rh'$ for the home, the work, and the energy bill.
A buyer pays the energy bill of the new home in the year of purchase and no
maintenance on it until the next year.

\emph{Staying.} An owner who stays pays
\begin{equation*}
  c+b'+\tilde P+\delta_hp^{t_0}h+e^{t_0}h=x,
  \qquad
  P(m,n)\leq\tilde P\leq(1+r_m)m,
\end{equation*}
and its balance next year is $m'=(1+r_m)m-\tilde P$. An owner who
renovates in place also pays $C(h)$, and the lower limit on its payment
falls to $P(m,n)-A$, where $A$ is the renovation advance of
Appendix~\ref{app:model-mortgages}.
In the year of the renovation it pays maintenance and the energy bill of
an inefficient home, and it receives the services of an efficient home.
